%%%%%%%%%%%%%%%%%%%%%%%%%%%%%%%%%%%%%%%%%%%%%%%%%%%
%
%  CHAPTER VI -- HYPERGRAPH ROTATION SYSTEMS
%
%%%%%%%%%%%%%%%%%%%%%%%%%%%%%%%%%%%%%%%%%%%%%%%%%%%

\chapter{\uppercase{Hypergraph Rotation Systems: A Unified Representation}}
\label{cha:hgrs}

This chapter presents the central contribution of the dissertation, Hypergraph Rotation Systems (H-GRS)~\cite{yildiz2026hypergraph}. An H-GRS replaces the family-specific replacement rules of Chapters~\ref{cha:modular-planar} to~\ref{cha:linked} with a single local operator. The only input is a graph whose faces are given, together with one integer twist per edge. Strands are generated by tracing the boundaries of the faces around each edge, and whether the result is a knot, a link, a weave, a knit, or a chainmail depends only on the twists.

\section{From Rotation Systems to Hypergraph Rotation Systems}

A graph does not determine its cycles. Tait already used graphs to describe knots~\cite{tait1877knots}, but the cycles of a graph become well defined only once a rotation system fixes the cyclic order of edges around each vertex~\cite{edmonds1960combinatorial,gross2001topological}. Rotation systems are what mesh data structures store, and adding a twist to each edge of a mesh turns its face boundaries into interlaced strands. This is how plain, twill, and single-cycle weaves were constructed on surfaces~\cite{akleman2009cyclic,xing2010single,akleman2015extended,akleman2020topologically}.

Those constructions have two restrictions. Each edge is shared by exactly two faces, and twists are binary or limited to $\{-1,0,+1\}$. Both come from deriving the chirality of a twist from the two sides of a surface. In three dimensions, however, a twist has an intrinsic chirality: a local reference direction along an edge is enough to tell a left-handed twist from a right-handed one. A global surface, and the two-faces-per-edge rule that comes with it, is therefore unnecessary. An H-GRS keeps the determinism of rotation systems while allowing any number of faces around an edge and any integer twist.

\section{Definition}
\label{sec:h-grs}

\begin{definition}[Hypergraph rotation system~\cite{yildiz2026hypergraph}]
A hypergraph rotation system is an ordered incidence structure $M=(F,E,\{\sigma_{\vec e}\},\{\rho_f\})$, where $F$ is a set of faces, $E$ a set of edges, $\sigma_{\vec e}$ a cyclic ordering of the faces incident to edge $e$ relative to a reference direction $\vec e$ (with $\sigma_{\overleftarrow{e}}=\sigma_{\vec e}^{-1}$), and $\rho_f$ an unoriented cyclic ordering of the edges bounding face $f$. The number of faces at an edge is its valence $m_e$.
\end{definition}

A \emph{knotted surface representation} is a pair $(M,t)$ with $t:E\to\mathbb{Z}$. At an edge with ordering $\sigma_e=(f_0,\ldots,f_{m-1})$, the twist acts as the cyclic shift $f_i\mapsto f_{i+t(e) \bmod m}$, reconnecting the face-sides that meet there. Tracing these reconnections across all edges concatenates face boundaries into closed strands. With $t\equiv0$ every face is its own cycle and the result is the scaffold itself. Nonzero twists merge, split, or reorder cycles, so the topology of the result is not stored in the scaffold but induced by the twists.

\begin{theorem}[Local cycle structure~\cite{yildiz2026hypergraph}]
At an edge of valence $m$ with twist $k$, the local permutation has exactly $\gcd(m,k)$ cycles, each of length $m/\gcd(m,k)$. If $\gcd(m,k)=1$, all incident faces merge into one cycle. If $k\equiv0 \pmod m$, every face remains separate.
\end{theorem}

Figure~\ref{fig:hgrs-twists} shows this rule at single edges. On an edge shared by two faces, even twists give two interlocked rings and odd twists give one knotted curve. On an edge shared by three faces, multiples of three give three interlocked rings and every other twist gives a single cycle.

\begin{figure}[htbp]
\centering
\includegraphics[width=6in]{ch06_fig01_twist_rule.png}
\caption[The twist rule at a single edge.]{The twist rule at a single edge. Top: an edge shared by two faces, twists $-2$ to $2$. Bottom: an edge shared by three faces, twists $-3$ to $3$. The number of resulting strands is $\gcd(m,k)$.
\figtodo{\cite{yildiz2026hypergraph} Fig.~\fs{fig:2_mantwist} (\fs{images/2man/input_mesh.jpg}, \fs{n2_tw.jpg}, \fs{n1_tw.jpg}, \fs{0_tw.jpg}, \fs{1_tw.jpg}, \fs{2_tw.jpg}) and Fig.~\fs{fig:3_mantwist} (\fs{images/3man/n3_tw.jpg} to \fs{3_tw.jpg}).}}
\label{fig:hgrs-twists}
\end{figure}

\section{Constructions}
\label{sec:hgrs-constructions}

\textbf{Chainmail.} If every twist is a multiple of its edge valence, $t(e)=a_e m_e$, no edge merges faces and every face remains an independent cycle. Nonzero multiples still change the embedding, so neighboring cycles interlock. Chainmail is therefore the set of assignments that preserve the cycle decomposition while changing how cycles pass through each other. A regular tetrahedron with twists of $\pm2$ already yields 12 distinct chainmail configurations.

\textbf{Unicursal knots.} The opposite goal is a single strand. When every edge has valence one or two, assigning odd twists to the edges of a spanning tree of the face-adjacency graph and even twists elsewhere merges all faces into one cycle. For higher valence, a twist can connect and disconnect cycles at the same time, so we split a problematic edge into two overlapping edges, each taking a contiguous part of its cyclic ordering. Splitting preserves the geometry and reduces valence, so the procedure always terminates. Every connected H-GRS therefore admits a unicursal realization. The tetrahedron yields 32 distinct unicursal knots.

\textbf{Design space.} With $k$ allowed twist values on $|E|$ edges and orientation-preserving symmetry group $G$, Burnside's lemma gives
\begin{equation}
N_G(k)=\frac{1}{|G|}\sum_{g\in G}k^{c_E(g)},
\end{equation}
where $c_E(g)$ is the number of edge cycles of $g$. This is the same counting used for cP-weaves in Chapter~\ref{cha:volumetric-bravais}, now applied to twist labels. With three twist values there are 87 distinct assignments on a tetrahedron, 22{,}815 on a cube, and about $3.43\times10^{12}$ on a dodecahedron.

\begin{figure}[htbp]
\centering
\includegraphics[width=6in]{ch06_fig02_tetra_chainmail_knots.png}
\caption[Chainmail and unicursal knots from a tetrahedron.]{One scaffold, different twists. Top: chainmail from a tetrahedron with twists of $\pm2$. Bottom: unicursal knots from the same tetrahedron with twists of $\pm1$ and $\pm2$.
\figtodo{\cite{yildiz2026hypergraph} Fig.~\fs{fig:tetra_link} (\fs{images/tetchain/0.png}--\fs{3.png}) and Fig.~\fs{fig:tetra_knot} (\fs{images/tetknot/0.png}--\fs{3.png}).}}
\label{fig:tetra-examples}
\end{figure}

\section{Realization and Behavior}
\label{sec:hgrs-behavior}

Strands are realized geometrically from a corner connectivity graph: each face-vertex incidence becomes a corner, the twists decide which corners connect across each edge, and strands wind helically around the edges between corners. The combinatorial guarantees concern connectivity; the geometric embedding is a starting point that may need further processing for a given application.

Because cycle structure is controlled locally, mechanical behavior can be designed as well. We decompose a scaffold into regions that are each unicursal and connect the regions through edges with twists that are multiples of their valence. Unicursal regions act as compliant units, and the connecting edges act as hinges. The graph of regions and hinges then predicts the behavior class: a tree gives open articulated chains, a surface gives origami-like panels, and a graph with loops gives flexagon-like mechanisms. Printed cubes with the same number of cycles but different linkage patterns behaved differently, one moving freely and one acting as a spring (Figure~\ref{fig:hgrs-prototypes}).

\section{Periodic H-GRS: Woven and Knitted Materials}
\label{sec:hgrs-periodic}

When the scaffold is a triply periodic honeycomb, twists are specified on a fundamental domain and repeated by translation, and strands run through the whole structure. With a uniform twist on the cubic honeycomb, where each edge is shared by four cells, the number of strand components changes systematically with the single twist value. The same uniform construction on any doubly or triply periodic honeycomb yields planar or volumetric plain woven fabrics. Non-uniform twists on distinct edge classes produce knits, chainmail, and layered structures that link across layers (Figure~\ref{fig:hgrs-periodic}).

In this setting the distinction between weaving and knitting is a property of the cycle structure. Multi-component assignments give interlaced systems of separate strands, as in weaving and chainmail. Single-component assignments give one continuous filament, as in knitting. Intermediate assignments give hybrids. The same periodic graph can also carry several different H-GRS structures, for example through different choices of face cycles on the Laves graph, so the map from periodic graphs to fibrillar materials is many-to-many.

\begin{figure}[htbp]
\centering
\includegraphics[width=6in]{ch06_fig03_periodic_hgrs.png}
\caption[Periodic H-GRS structures.]{Periodic H-GRS. Top: uniform twists on the cubic honeycomb. Middle: non-uniform twists on a 2D Wigner-Seitz tiling give knit-like and chainmail structures. Bottom: non-uniform twists on the cubic honeycomb give single-component and multi-component volumetric structures.
\figtodo{\cite{yildiz2026hypergraph} Fig.~\fs{fig:bravais3D_repeated} (\fs{images/bravais/cp_1.jpg}--\fs{cp_3.jpg}), Fig.~\fs{fig:bravais2D_repeated_nonuniform} (\fs{images/ws_2d/2_1.jpg}, \fs{2_2_1.jpg}), Fig.~\fs{fig:cp_single} (\fs{images/cp_nonuniform/2_1_1_tiled.jpg}) and Fig.~\fs{fig:cp_multiple} (\fs{images/cp_nonuniform/1_3_1_tiled.jpg}).}}
\label{fig:hgrs-periodic}
\end{figure}

\section{Synthesis: H-GRS as the Unifying Representation}
\label{sec:hgrs-synthesis}

We now return to Table~\ref{tab:scheme}. Table~\ref{tab:hgrs-synthesis} lists, for each earlier construction, what it corresponds to in an H-GRS and how complete that correspondence is.

\begin{table}[htbp]
\centering
\begin{singlespace}
\footnotesize
\begin{tabular}{@{}>{\raggedright\arraybackslash}p{0.2\textwidth}>{\raggedright\arraybackslash}p{0.47\textwidth}>{\raggedright\arraybackslash}p{0.22\textwidth}@{}}
\toprule
\textbf{Construction} & \textbf{In H-GRS terms} & \textbf{Correspondence} \\
\midrule
Planar and surface weaves (Ch.~\ref{cha:modular-planar}) & Scaffold with every edge of valence two; twists in $\{0,\pm1\}$. Classical cyclic weaving is this special case. & Contained \\[3pt]
Volumetric weaves (Ch.~\ref{cha:volumetric-bravais}) & Periodic H-GRS on a 3-honeycomb. Uniform twists give volumetric plain weaves; non-uniform twists give further woven families. & Conceptual; whether every cP-weave arises is open \\[3pt]
Labeled lattices (Ch.~\ref{cha:nonwoven}) & Twists on a fundamental domain form a periodic label array, so the $P^K$-adic operations apply to them. & Observation \\[3pt]
Linked structures (Sec.~\ref{sec:linked}) & Twists that are nonzero multiples of valence interlock face-cycles directly; zero twists play the role of unbundled edges. & Analogous; no separate edge-cycles \\[3pt]
Knots and knits & Unicursal assignments, finite or periodic. & Contained \\[3pt]
Nonwovens (Ch.~\ref{cha:nonwoven}) & Attachments rather than crossings; frames rather than cycles. & Outside \\[3pt]
Polycatenanes (Sec.~\ref{sec:polycatenanes}) & Components are spatial graphs, not cycles. & Outside \\
\bottomrule
\end{tabular}
\end{singlespace}
\caption[Earlier constructions expressed as hypergraph rotation systems.]{Earlier constructions expressed as hypergraph rotation systems. ``Contained'' means the construction is a special case; ``conceptual'' and ``analogous'' mean the same scheme applies but the structures produced differ in detail.}
\label{tab:hgrs-synthesis}
\end{table}

Three points follow. First, every cycle-based family in this dissertation, which covers weaves, knits, links, knots, and chainmail, is produced by the same operator on the same kind of input. What used to be a choice of replacement rule is now a choice of twist values. Second, the machinery of the earlier chapters carries over: Wigner-Seitz honeycombs serve as scaffolds, twist labels on a fundamental domain are a labeled array, and Burnside counting applies to twist assignments exactly as it applied to cP configurations. Third, the boundary of the unification is clear. Nonwovens and polycatenanes are not cycle structures, and the paper itself notes that cycle-based descriptions do not capture every form of entanglement. We therefore describe H-GRS as a unification of the cycle-based constructions and a common conceptual frame for the rest, not as a replacement for every earlier result.

\section{Physical Prototypes}

We built H-GRS structures as single 3D printed pieces, as unfolded parts closed after printing, as modular assemblies, and as paper strips over a paper scaffold of triangular prisms that follows the valence-three edges of the truncated octahedral honeycomb (Figure~\ref{fig:hgrs-prototypes}). These prototypes confirmed that the combinatorial predictions hold physically, including the difference in behavior between assignments with the same number of cycles. They also showed the limits of ad hoc fabrication. The framework says nothing about thickness, clearance, or support, so each prototype needed its own manufacturing decisions. Turning these decisions into methods is the subject of the proposed work.

\begin{figure}[htbp]
\centering
\includegraphics[width=6in]{ch06_fig04_prototypes.png}
\caption[Physical H-GRS prototypes.]{Physical H-GRS prototypes. Left: printed cubes with different twist assignments, freely moving and compliant. Middle: a printed articulated structure from a flexagon scaffold. Right: paper-strip realizations of the truncated octahedral honeycomb with one, two, and three twists.
\figtodo{\cite{yildiz2026hypergraph} Fig.~\fs{fig:cube_behaviour} (\fs{images/photos/cube_1.jpg}, \fs{cube_2.jpg}), Fig.~\fs{fig:flexon_behaviour} (\fs{images/photos/flex_1.jpg}) and Fig.~\fs{fig:troct_paper} (\fs{images/photos/1twist_clear.jpeg}, \fs{2twist_cleaner.jpeg}, \fs{3twist_cleaner.jpeg}).}}
\label{fig:hgrs-prototypes}
\end{figure}
